\documentclass[11pt]{article}
\usepackage[a4paper,margin=1in]{geometry}
\usepackage[T1]{fontenc}
\usepackage{lmodern,microtype}
\usepackage{amsmath,amssymb,amsthm,mathtools}
\usepackage{booktabs,enumitem,xcolor}
\usepackage[numbers,sort&compress]{natbib}
\usepackage[colorlinks=true,linkcolor=blue!55!black,citecolor=blue!55!black,urlcolor=blue!55!black]{hyperref}
\linespread{1.06}
\newtheorem{theorem}{Theorem}[section]
\newtheorem{proposition}[theorem]{Proposition}
\newtheorem{corollary}[theorem]{Corollary}
\theoremstyle{remark}\newtheorem{remark}[theorem]{Remark}
\newcommand{\C}{\mathbb C}

\title{Tight Clouds versus Locally Finite Divisors\\
\large An Exact Obstruction to Direct Rank-Growing Root Limits}
\author{Liang Wang}
\date{July 2026}

\begin{document}
\maketitle

\begin{abstract}
Global centered-RMS normalization makes every empirical root measure have
unit second moment, and RH-224 therefore proves uniform tightness. This paper
shows why that positive probability result cannot be promoted directly to a
holomorphic zero-divisor limit.

Let $\Lambda_n$ be finite clouds with ranks $d_n\to\infty$, let
$\nu_n=\sum_{\lambda\in\Lambda_n}\delta_\lambda$ be their integer root
divisors, and let $\mu_n=d_n^{-1}\nu_n$ be their empirical probability
measures. If $\{\mu_n\}$ is uniformly tight, one fixed compact set contains at
least half of every cloud. Hence $\nu_n$ has mass at least $d_n/2$ on a fixed
compact set, which diverges. The divisors cannot converge vaguely to a locally
finite divisor, and they cannot be zero divisors of a locally uniformly
convergent nonzero holomorphic family.

In the RH-222 atlas all normalized roots lie in $|z|\le2$, while all raw
Hardy-scaled resonances lie in $|z|<1$ with maximum modulus $0.86860$.
Compact divisor counts grow from $4$ to $35$ on the left and from $4$ to $34$
on the right.

The obstruction rejects the resonances themselves, normalized or raw, as the
zeros of the sought growing determinant. It does not reject the reciprocal
Fredholm zeros $1/\lambda$, which can escape to infinity when resonances tend
to zero. This distinction redirects the next layer to the correct spectral
variable.
\end{abstract}

\section{Two measures carried by the same roots}

For a finite multiset
\[
 \Lambda_n=\{\lambda_{n,1},\ldots,\lambda_{n,d_n}\}\subset\C,
\]
there are two natural measures:
\begin{equation}\label{eq:measures}
 \nu_n=\sum_{j=1}^{d_n}\delta_{\lambda_{n,j}},\qquad
 \mu_n=\frac1{d_n}\nu_n.
\end{equation}
The first is the integer divisor measure. The second is the empirical
probability measure.

Weak compactness of $\mu_n$ is a macroscopic statement in which each root has
mass $1/d_n$. Local finiteness of $\nu_n$ is a microscopic statement in which
each simple root has mass one. RH-224 proves tightness only for the first
normalization \citep{WangRH224}.

\section{The tightness obstruction}

\begin{theorem}[Tight rank growth is not a locally finite divisor]
\label{thm:obstruction}
Suppose $d_n\to\infty$ and the empirical measures $\{\mu_n\}$ in
\eqref{eq:measures} are uniformly tight. Then there is a compact set
$K\subset\C$ such that
\begin{equation}\label{eq:massgrowth}
 \nu_n(K)\ge\frac{d_n}{2}
\end{equation}
for every $n$. Consequently no subsequence of $\nu_n$ converges vaguely to a
locally finite Radon measure.
\end{theorem}
\begin{proof}
Uniform tightness with $\varepsilon=1/2$ supplies compact $K$ such that
$\mu_n(K)\ge1/2$ for all $n$. Multiplying by $d_n$ proves
\eqref{eq:massgrowth}.

Choose a compactly supported continuous function $\phi\ge0$ that equals one
on $K$. Then
\[
 \int\phi\,d\nu_n\ge\nu_n(K)\longrightarrow\infty.
\]
Vague convergence to a locally finite Radon measure $\nu$ would instead give
the finite limit $\int\phi\,d\nu$. This is impossible.
\end{proof}

The fraction $1/2$ is arbitrary. For any fixed
$0<\varepsilon<1$, one compact set carries at least
$(1-\varepsilon)d_n$ divisor mass.

\begin{corollary}[No direct nonzero holomorphic limit]\label{cor:holo}
Under the assumptions of Theorem~\ref{thm:obstruction}, the $\Lambda_n$
cannot be the complete zero divisors on $\C$ of holomorphic functions
$f_n$ converging locally uniformly to a nonzero entire function $f$, with
zeros counted by algebraic multiplicity.
\end{corollary}
\begin{proof}
On a compact neighborhood whose boundary avoids the zeros of $f$, local
uniform convergence and Rouche's theorem make the number of zeros eventually
equal to that of $f$; see \citet{Conway1978}. Equivalently, zero divisors of a
nonzero locally uniform limit are locally bounded. This contradicts
Theorem~\ref{thm:obstruction}.
\end{proof}

The only holomorphic escape is degeneration to the identically zero function,
which is not a spectral determinant.

\section{Application to globally normalized clouds}

RH-224 gives
\[
 \mu_n(|z|>R)\le R^{-2}.
\]
At $R=2$, every normalized cloud has at least $75\%$ of its roots in one
fixed disk. Thus any all-level extension with ranks tending to infinity would
satisfy
\[
 \nu_n(\overline{D(0,2)})\ge\frac34d_n\to\infty.
\]
The exact theorem already rejects the direct normalized-root divisor. No
numerical convergence assumption is needed.

This conclusion complements the fixed-quartic obstruction RH-219
\citep{WangRH219}. RH-219 rules out degree four and repeated fixed support.
Theorem~\ref{thm:obstruction} rules out a different failure mode: genuinely
new roots are added, but a positive fraction remains in one compact region.

\section{Finite physical audit}

The sixteen-level clouds have strictly increasing ranks. The normalized and
raw compact counts are:
\begin{center}
\begin{tabular}{lrrr}
\toprule
channel & first rank & final rank & compact-count growth\\
\midrule
left, normalized $|z|\le2$ & 4 & 35 & 31\\
right, normalized $|z|\le2$ & 4 & 34 & 30\\
left, raw $|z|\le1$ & 4 & 35 & 31\\
right, raw $|z|\le1$ & 4 & 34 & 30\\
\bottomrule
\end{tabular}
\end{center}
Every audited normalized root lies in $|z|\le2$. Every raw selected resonance
lies in the unit disk, and the largest raw modulus is $0.868592$.

The raw statement is finite. Although Markov resonances before Hardy scaling
lie in the unit disk at fixed positive noise, the chosen $0.85$ scaling could
in principle move some bulk roots beyond one at smaller unobserved scales.
The theorem-level direct-normalization obstruction does not depend on that
finite raw-radius observation.

\section{Why reciprocal zeros are different}

For a compact operator with nonzero eigenvalues $\lambda_j\to0$, the Fredholm
determinant has zeros at
\begin{equation}
 z_j=\lambda_j^{-1}.
\end{equation}
The newly added reciprocal zeros can escape every compact set. Their integer
divisor may therefore remain locally finite even though the resonance
empirical measure is tight or even converges to a point at zero.

Inversion is singular at the resonance accumulation point. It does not
preserve probability tightness, and Theorem~\ref{thm:obstruction} cannot be
applied to the reciprocal cloud without a separate hypothesis.

This is precisely the fixed-noise architecture of the regularized determinant
constructed in RH-7 \citep{WangRH7}. The next paper restores the finite
identity between the characteristic polynomial in $\lambda$ and the
Fredholm polynomial in $z$.

\section{Route consequence}

The following routes are now separated:
\begin{enumerate}[leftmargin=2em]
\item raw or centered resonances used directly as zeros: rejected;
\item repeated fixed factors: rejected by RH-219;
\item reciprocal resonance zeros at fixed noise: operator-theoretically valid;
\item a reciprocal small-noise divisor: still open and must pass local-count
and omitted-tail tests.
\end{enumerate}

No small-noise obstruction for reciprocal zeros is proved here. Gate A
remains open, as do all self-adjoint, counting-law, arithmetic, and zeta
identification stages.

\bibliographystyle{plainnat}
\bibliography{references}
\end{document}
